%
% !TEX root = Main.tex 
\section{Change of measure}
\label{app:changeMeasureLem}
%
The next lemma is a finite-horizon version of the observed-history
change-of-measure method; see
\citet[Lemma~1 and Appendix~A]{Kaufmann16CBI}.
We give a self-contained proof allowing common finite-support reward
laws and include the Bretagnolle--Huber testing bound; see
\citet[Theorem~14.2]{Lattimore20BA}.

\begin{lemma}[\textcolor{titleTh}{Change of measure for observed transcripts}]
\label{l:transcriptChangeMeas}
Fix a positive integer horizon $\timeHorizon$, $\nArms\geq2$ arms,
an arm $k\in\setArms$, and a deterministic set of rounds
$\mathcal T\subseteq\{1,\ldots,\timeHorizon\}$, possibly the entire
horizon. The arm and set are fixed before interaction. In each of two
environments $A$ and $B$, the reward-table entries are independent,
independent of the learner's internal randomness, and have
finite-support distributions in $[0,1]$. At the pairs $(k,t)$ with
$t\in\mathcal T$, environment $A$ uses Bernoulli$(u)$ and environment
$B$ uses Bernoulli$(v)$, with fixed $u,v\in(0,1)$. At every other
arm-round pair, the two distributions are identical; these common
distributions may depend on the arm and round and may be deterministic.
Since there are finitely many entries, a finite reward alphabet
contains all their supports in both environments, including any
values other than $0$ and $1$.

Run the same adaptive randomized learner with the same horizon in both
environments. Its actions use only past actions, observed rewards,
and internal randomness, which may persist across rounds. It is not
given the environment's identity. Its final recommendation
$\recomArm_{\timeHorizon}\in\setArms$ also uses only the observed
interaction and internal randomness.
Let $\sample_t=\gainVector_{\pulledArm_t,t}$ be its observed reward and set
\[
Z=(\pulledArm_1,\sample_1,\ldots,
\pulledArm_{\timeHorizon},\sample_{\timeHorizon},\recomArm_{\timeHorizon}).
\]
Only observed rewards belong to $Z$, not the complete table.
Let $P_A$ and $P_B$ be its laws on the finite transcript space
$\mathcal Z$, and let $\Exp_A$ denote expectation in environment $A$.
Define
\[
N_k(\mathcal T)=\sum_{t\in\mathcal T}\indicator{\pulledArm_t=k}.
\]
Use natural logarithms and binary relative entropy
\[
\operatorname{kl}(a,b)
=a\log\frac{a}{b}+(1-a)\log\frac{1-a}{1-b},
\qquad a\in[0,1],\quad b\in(0,1),
\]
where zero-weight terms are zero. For laws with common support on
$\mathcal Z$, define
\[
D(P_A\Vert P_B)
=\sum_{z:P_A(z)>0}P_A(z)\log\frac{P_A(z)}{P_B(z)}.
\]
Then:
\begin{enumerate}[label=(\roman*)]
\item $P_A$ and $P_B$ have the same support.
\item Their relative entropy satisfies
\[
D(P_A\Vert P_B)=\Exp_A[N_k(\mathcal T)]\operatorname{kl}(u,v).
\]
\item For every event $W$ determined by the observed transcript, set
$p=P_A(W)$ and $q=P_B(W)$. If $0<q<1$, then
\[
\operatorname{kl}(p,q)\leq D(P_A\Vert P_B),
\qquad
p\log(1/q)\leq\Exp_A[N_k(\mathcal T)]\operatorname{kl}(u,v)+\log2.
\]
If $q=0$, then $p=0$; if $q=1$, then $p=1$.
\item For every observed-transcript event $W$, the
Bretagnolle--Huber testing bound holds:
\begin{equation}\label{eq:transcriptTesting}
P_A(W^c)+P_B(W)\geq\frac12\exp\left(-D(P_A\Vert P_B)\right).
\end{equation}
\end{enumerate}
The event must be an observed-transcript event; this statement does
not directly apply to events involving unobserved table entries.
\end{lemma}

\begin{proof}
Draw all learner randomness as a single seed $R$ before interaction,
independently of the reward table. Once this seed is fixed, actions,
internal state updates, and the final recommendation are determined
by the observed history. This includes randomness retained across
rounds and final randomization. For a fixed
transcript $z=(i_1,x_1,\ldots,i_{\timeHorizon},x_{\timeHorizon},j)$,
let $C_z$ be the set of seeds for which the learner produces the
recorded actions and recommendation when sequentially supplied the
recorded observations. Set $w(z)=\Pro_R(C_z)$ and let
$f^E_{j,t}(x)$ be the probability mass function of entry $(j,t)$
in environment $E$ on the finite reward alphabet.
For each seed in $C_z$, the observed table coordinates $(i_t,t)$
are distinct and independent. Their probability of matching the
recorded rewards is the product of their probability masses;
all unobserved entries integrate to one. Integrating over these seeds
therefore gives, for $E\in\{A,B\}$,
\[
P_E(z)=w(z)\prod_{t=1}^{\timeHorizon}
f^E_{i_t,t}(x_t).
\]
The factor $w(z)$ is the same in both environments and need not sum
to one over transcripts. At each arm-round pair, the two reward
distributions have the same support: they are either identical or
Bernoulli with parameters in $(0,1)$. Although common reward factors
can be zero, the factorization therefore implies
\[
P_A(z)>0\quad\Longleftrightarrow\quad P_B(z)>0.
\]
This proves common support without requiring the learner to assign
positive probability to every action.

On this common positive support, $w(z)$ and all observed reward
factors are positive. The learner factor and identical reward factors
cancel from the likelihood ratio, leaving only observed pulls of arm
$k$ at rounds in $\mathcal T$. Thus
\[
\log\frac{P_A(Z)}{P_B(Z)}
=\sum_{t\in\mathcal T}\indicator{\pulledArm_t=k}
\left[\sample_t\log\frac{u}{v}
+(1-\sample_t)\log\frac{1-u}{1-v}\right].
\]
At each round, condition on the seed, the past interaction, and the
selected arm. Under $A$, if $t\in\mathcal T$ and $\pulledArm_t=k$,
the current observed reward is Bernoulli$(u)$, since its table entry
is independent of the seed and past. Taking expectations yields
\[
D(P_A\Vert P_B)
=\Exp_A\left[\sum_{t\in\mathcal T}
\indicator{\pulledArm_t=k}\operatorname{kl}(u,v)\right]
=\Exp_A[N_k(\mathcal T)]\operatorname{kl}(u,v).
\]
Common deterministic rounds contribute zero additional information,
even though the learner's actions and final recommendation may
depend on their rewards.

For $0<q<1$, regard $W$ as a subset of $\mathcal Z$ and let
$\mathcal S$ be the common positive support. Applying the finite
log-sum inequality separately to $W$ and its complement gives
\begin{align*}
D(P_A\Vert P_B)
&=\sum_{z\in W\cap\mathcal S}P_A(z)\log\frac{P_A(z)}{P_B(z)}
+\sum_{z\in\mathcal S\setminus W}P_A(z)\log\frac{P_A(z)}{P_B(z)}\\
&\geq p\log\frac{p}{q}+(1-p)\log\frac{1-p}{1-q}
=\operatorname{kl}(p,q).
\end{align*}
Zero-mass groups are interpreted by continuity. Moreover,
\begin{align*}
\operatorname{kl}(p,q)
&=p\log(1/q)+(1-p)\log(1/(1-q))\\
&\quad+p\log p+(1-p)\log(1-p)\\
&\geq p\log(1/q)-\log2.
\end{align*}
The second term in the first line is nonnegative, and binary entropy
is maximized at $p=1/2$ with value $\log2$. Endpoint values of $p$ use the
zero-weight convention. This proves the event consequence.
When $q=0$ or $q=1$, common support gives respectively $p=0$ or $p=1$
without evaluating undefined logarithms.

For the testing bound, on the common positive support $\mathcal S$,
define
\[
a=\sum_{z\in\mathcal S}\sqrt{P_A(z)P_B(z)},\qquad
s=\sum_{z\in\mathcal S}\min\{P_A(z),P_B(z)\}.
\]
Jensen's inequality gives
\[
\log a
=\log\Exp_A\left[\sqrt{\frac{P_B(Z)}{P_A(Z)}}\right]
\geq-\frac12D(P_A\Vert P_B),
\]
so $a^2\geq\exp(-D(P_A\Vert P_B))$. By Cauchy--Schwarz,
\[
a^2\leq
\left(\sum_{z\in\mathcal S}\min\{P_A(z),P_B(z)\}\right)
\left(\sum_{z\in\mathcal S}\max\{P_A(z),P_B(z)\}\right)
=s(2-s)\leq2s.
\]
Finally, $P_A(W^c)+P_B(W)\geq s$, proving
Equation~\ref{eq:transcriptTesting}. This also covers $W$ empty or
equal to the whole transcript space without taking logarithms of
event probabilities.

The distinguishing information equals the expected number of
informative pulls times the information supplied by each pull.
For $\mathcal T=\{1,\ldots,m\}$, $N_k(\mathcal T)=\pullsNumber_k(m)$.
\end{proof}
